\chapter{完整源码清单：Lab 00 可信 Benchmark}

本附录完整收录 \filepath{books/cpu-volume-1/labs/lab00_benchmark_foundation}。第一册的主线不是“写几个 kernel”，而是从 C++ 源码一路走到可复查的机器证据：函数合同、checksum、runner、CSV、坏 benchmark、测试和运行脚本都要同时出现。

\section{一、CMakeLists.txt}

构建文件组织库、bench 程序、坏 benchmark 程序和测试目标。读者要看清 Debug/Release 构建不会改变源码合同，但会改变机器码和性能证据。

\lstinputlisting[language=bash,caption={lab00\_benchmark\_foundation/CMakeLists.txt}]{../../labs/lab00_benchmark_foundation/CMakeLists.txt}

\section{二、README}

README 给出实验目的、运行命令和报告口径。它不是装饰文档，而是读者复现 CSV 和反汇编证据的入口。

\lstinputlisting[language=bash,caption={lab00\_benchmark\_foundation/README.md}]{../../labs/lab00_benchmark_foundation/README.md}

\section{三、Kernel 公共头}

\filepath{kernels.hpp} 定义最小计算内核。所有函数都用 \code{std::span} 表达连续内存视图，避免复制大数组；读者应从这里确认每个 kernel 的输入、输出和错误边界。

\lstinputlisting[language=C++,caption={include/lab00/kernels.hpp}]{../../labs/lab00_benchmark_foundation/include/lab00/kernels.hpp}

\section{四、Benchmark 公共头}

\filepath{benchmark.hpp} 定义 runner、case、result 和 CSV 输出边界。它把“测量多少次、记录什么字段、checksum 如何保留”变成代码合同。

\lstinputlisting[language=C++,caption={include/lab00/benchmark.hpp}]{../../labs/lab00_benchmark_foundation/include/lab00/benchmark.hpp}

\section{五、Kernel 实现}

这里是求和、点积、填充、复制和 checksum 的完整实现。读者应重点观察：哪些循环是被测热路径，哪些计算结果被 checksum 保留，哪些错误会直接抛出异常。

\lstinputlisting[language=C++,caption={src/kernels.cpp}]{../../labs/lab00_benchmark_foundation/src/kernels.cpp}

\section{六、Benchmark Runner 实现}

runner 实现负责预热、正式测量、过滤 case、记录耗时和写 CSV。它不应该知道某个 kernel 的内部算法；它只管理测量边界。

\lstinputlisting[language=C++,caption={src/benchmark.cpp}]{../../labs/lab00_benchmark_foundation/src/benchmark.cpp}

\section{七、主 Benchmark 程序}

主程序把数据准备放在计时区域之外，并注册可比较的 workload。读者应检查命令行参数如何进入 runner，以及 checksum 如何随每次测量输出。

\lstinputlisting[language=C++,caption={src/main.cpp}]{../../labs/lab00_benchmark_foundation/src/main.cpp}

\section{八、坏 Benchmark 程序}

坏 benchmark 故意把无关成本放进计时区域，用来训练读者识别错误测量。它能产生结果，但结果不能直接解释 kernel 性能。

\lstinputlisting[language=C++,caption={src/bad\_benchmarks.cpp}]{../../labs/lab00_benchmark_foundation/src/bad_benchmarks.cpp}

\section{九、完整测试}

测试文件覆盖 kernel 正确性、错误长度、checksum、runner 迭代次数、过滤和 CSV 格式。没有这些测试，benchmark 数字不能当作证据。

\lstinputlisting[language=C++,caption={tests/lab00\_tests.cpp}]{../../labs/lab00_benchmark_foundation/tests/lab00_tests.cpp}

\section{十、运行脚本}

脚本把构建、测试、运行和结果目录串起来。它适合作为本地 smoke，但正式报告仍要记录 CPU、编译器、构建类型和环境噪声。

\lstinputlisting[language=bash,caption={scripts/run\_lab00.sh}]{../../labs/lab00_benchmark_foundation/scripts/run_lab00.sh}
